\section{一、摘要}

黑土地是保障国家粮食安全的重要战略资源，东北黑土地作为我国最重要的商品粮基地之一，其土层厚度、有机质含量与结构稳定性的变化，直接关系到区域农业生产的可持续性和国家粮食供给的长期稳定。近年来，随着《中华人民共和国黑土地保护法》的颁布实施，以及《国家黑土地保护工程实施方案（2021—2025年）》《东北黑土地保护性耕作行动计划（2020—2025年）》等系列政策的持续推进，黑土地保护已由局部试点上升为制度化、常态化的国家任务，对监测手段的时效性、覆盖范围与数据可信度提出了更高要求。然而，传统黑土地监测长期依赖人工实地巡查与目视解译，存在人工巡查效率低、光学遥感解译精度不足、数据与治理脱节等核心痛点：人工巡查受人力和交通条件制约，难以覆盖广袤耕地，且判读结果主观性强、复现性差；遥感数据获取能力虽不断增强，却缺乏与之匹配的自动化智能解译手段，致使海量影像“存而不用”，监测结论难以快速转化为可执行的治理决策。监测能力与治理需求之间的这一落差，正是本项目立项所要解决的关键问题。

从量化背景看，黑土地退化的累积效应已经显现：黑土层平均减薄6.83cm（累计值），有机质下降的样点占比达52.5\%，存在水土流失的样点占比达24.4\%。三组数据分别从土层厚度、肥力水平与表土流失三个维度刻画了退化的程度与普遍性，说明黑土地质量下降并非个别地块的偶发现象，而是区域尺度上的系统性问题。退化过程还具有渐变性与分散性：单次实地巡查难以察觉细微差异，而持续累积则可能演变为难以逆转的沟蚀与肥力衰退。因此治理必须建立在“早发现、早干预”的基础上，也要求监测具备周级频次与地块级分辨率，使保护措施能在退化尚处可逆阶段时及时介入。

针对上述痛点，本项目设计了一套基于人工智能的卫星遥感动态变化监测系统，构建“天空地一体化光学遥感采集+AI智能解译+可视化决策支撑”的完整技术闭环。其中，“天空地一体化采集”负责多平台、多时相影像的稳定获取；“AI智能解译”负责将原始影像转化为可量化的变化信息；“可视化决策支撑”负责把变化信息组织为面向基层管理者的直观图件与分析结论。三者相互衔接、逐级递进，形成从数据输入到决策输出的贯通链路。

在技术实现上，系统以国产GF-2亚米级光学遥感卫星影像与无人机航拍影像为核心数据源，以改进BIT双时相Transformer模型为算法核心，搭配ResNet18骨干网络实现高分辨率影像的多尺度特征提取；同时开发了前后端分离的轻量化可视化系统，支持批量影像处理、变化检测、结果可视化与一键生成分析报告，普通计算机即可完成本地化部署，降低了在基层单位推广使用的硬件门槛。经试点验证，系统对黑土地变化识别精确率达70\%，算力成本降低69\%，监测周期缩短33\%，已在黑龙江松嫩平原完成落地应用，可为黑土地常态化保护与精细化治理提供高效、自主可控的光电监测技术支撑，助力保障国家粮食安全。

\section{二、项目概况与核心组成}

\subsection{1.项目总体目标}

本项目面向黑土地保护“精准治理”的现实需求，以“天空地一体化监测+AI智能解译”为核心，开发一套土地光学遥感变化检测智能系统，实现黑土地退化的自动化识别、可视化分析与辅助决策，构建“数据采集—智能检测—分析决策”的完整技术闭环，为黑土地保护提供高效、精准的科技支撑。

上述目标可进一步分解为三个相互支撑的层面。第一是数据采集层面，打通卫星与无人机两类平台的影像获取通道，保证监测区域在时间维度上“有数据可用、有数据可比”；第二是智能检测层面，用深度学习模型替代重复性的人工目视解译，把侵蚀沟发育、植被覆盖变化、土地利用变化等退化信息从影像中自动提取出来，并使检测结果在精度与稳定性上满足业务要求；第三是分析决策层面，将模型输出转化为管理者可直接理解的图件、指标与报告，使监测结论能够进入日常巡查、执法核查与保护性耕作评估等具体工作流程。

需要说明的是，本项目并非面向通用场景的遥感监测平台，而是紧扣黑土地保护业务的专用系统，其差异化定位体现在三个层面。算法层面，通用变化检测模型多以城市扩张、地表覆盖转换等目标优化，本项目则针对侵蚀沟发育、植被覆盖变化、土地利用变化等黑土地退化信息专项训练，使模型关注的类别与业务关注的对象一致。数据层面，通用方案多单依赖卫星影像，受重访周期与云层遮挡制约，本项目以国产GF-2亚米级影像承担周期性普查、以无人机影像承担重点核查。结果层面，通用方案输出多为栅格差异图或矢量图斑，仍需二次判读，本项目则直接输出图件、指标与文字结论，把专业解译封装在系统内部。三方面取舍共同指向同一目标：使不具备遥感背景的基层使用者同样能获得可用的监测结论。

项目目标的设定强调“可用”与“可推广”并重。一方面，系统功能围绕真实业务环节组织，避免停留在算法演示层面；另一方面，系统在算力占用、部署方式与操作复杂度上做了针对性约束，使其能够在算力条件有限的基层单位稳定运行，从而具备从单点试点走向区域推广的现实基础。项目按照既定时间线推进：2022年启动政策调研与数据积累，2024年完成影像标注与模型验证，2026年获软件著作权并落地试点，2028年推广至黑吉辽蒙四省区，2030年构建“天空地”一体化监测网。

从成果基础看，项目并非从零起步。团队围绕遥感变化检测与监测系统开发已形成自主知识产权积累，累计获得软件著作权3项，并取得实用新型专利1项（专利号ZL 2025 2 2047856.0，授权公告号CN 224511490 U，专利权人为黑龙江八一农垦大学）。相关权利覆盖算法实现、系统软件与装置结构等层面，表明项目在原理验证与工程实现上均已具备可核验的基础，为后续推广提供了支撑。

\subsection{2.项目核心组成部分}

围绕上述目标，项目由AI算法核心、多源数据支撑、可视化智能系统三部分组成。三者分别对应“算得准”“看得清”“用得上”三个环节，共同构成完整的技术链条。

三者之间并非简单的功能并列，而是相互约束、彼此支撑的有机整体。算法对数据提出要求：Transformer模型需要在有限算力下完成全局建模，就要求影像在空间分辨率与时间一致性上达到一定标准，这直接决定了采集环节对GF-2亚米级影像与无人机核查影像的取舍；数据反过来检验算法：多时相、多平台的影像覆盖，使模型可能遇到的成像差异与季节变化被充分暴露，从而推动训练策略的持续调整。算法与数据共同支撑系统：只有模型输出稳定、数据基础可靠，可视化系统呈现的图件与结论才具备决策价值；而系统的实际使用又会反馈新的疑点与误差样本，为模型迭代提供标注来源。三者由此形成“数据喂养算法、算法支撑系统、系统反馈数据”的闭环，任何一环的削弱都会影响最终治理决策的可靠性。

\subsubsection{（1）AI算法核心：基于Transformer的深度学习模型}

以BIT（Bipartite Transformer）深度学习模型为算法核心，针对黑土地遥感影像的特征进行优化训练，实现对侵蚀沟、植被覆盖变化、土地利用变化等关键退化信息的高精度、自动化检测，解决传统人工解译效率低、主观性强的问题。

从技术原理看，变化检测的本质是对同一区域两个不同时相的影像进行逐像素的对应分析与差异判别。卷积神经网络虽擅长提取局部纹理与边缘特征，但其感受野受限于卷积核尺寸，对影像中大范围、长距离的空间关联刻画不足，容易在连片的耕地与沟道区域出现漏检或边界模糊。BIT双时相Transformer模型\cite{ref1}通过自注意力机制建模任意两个位置之间的依赖关系，能够在整幅影像范围内建立全局上下文联系，从而更准确地判断“变化区域是什么、变化边界在哪里”。模型中引入的语义分词器（semantic tokens）负责将影像特征归纳为一组高层语义表示，以此约束注意力计算的范围并保留类别语义信息；ResNet18骨干网络则负责从高分辨率影像中提取多尺度特征，兼顾细节纹理与宏观结构。二者的组合在保持Transformer全局建模能力的同时，抑制了计算量的膨胀，为在有限算力条件下完成推理提供了可行路径。

在训练环节，模型针对黑土地遥感影像的成像特点与地物分布规律进行专项优化：考虑不同季节植被覆盖差异、地表含水量变化以及云雾干扰等因素对影像色调的影响，通过标注样本的扩充与训练策略的调整，提升模型在复杂成像条件下的稳定性与泛化能力。与人工目视解译相比，模型不受主观经验差异影响，判读标准统一、结果可复现，且可批量处理连续时相影像，为周级动态监测提供了算法基础。

在方案取舍上，项目选择CNN骨干与Transformer注意力相结合，而非纯卷积或纯注意力的单一路线，原因在于黑土地遥感影像同时包含两类信息：地块边界、沟道边缘等依赖局部细节的纹理特征，以及连片耕地、沟系走向等依赖大范围空间关联的结构特征。纯卷积结构对后者刻画不足，纯注意力结构则在高分辨率影像上算力开销过大。以ResNet18骨干提取多尺度局部特征、由语义分词器汇聚高层语义、再交由双时相注意力完成差异判别，使模型在精度与算力之间取得平衡，这也是系统能够在普通计算机上完成推理、进而支撑本地化部署的前提。

\subsubsection{（2）多源数据支撑：“天空地”一体化光学遥感采集}

以高分辨率光学遥感影像为核心数据源，构建多平台协同采集体系：

\begin{itemize}
    \item 卫星端：支持GF-2等国产高分辨率卫星的批量数据处理，实现大范围周期性的黑土区动态监测；
    \item 无人机端：完成算法的机载适配，支持野外现场快速检测，弥补卫星数据重访周期长、受天气影响大的短板，实现重点区域的精细化核查。
\end{itemize}

两类平台在观测尺度与响应速度上形成互补。卫星平台观测范围大、成像稳定，适合承担区域性的周期性普查任务，其亚米级分辨率足以支撑耕地地块尺度上的变化识别；无人机平台部署灵活、受云层遮挡影响小，可在卫星影像存疑或天气条件不佳时快速抵达现场，对疑似变化图斑进行低空高分辨率复核。由此形成“卫星普查发现疑点、无人机核查确认结果”的作业模式，既保证了监测覆盖面，又提高了结论的确定性。同时，两类数据统一纳入同一套预处理与解译流程，避免因数据来源不同而产生的标准不一致问题，为后续算法处理与结果比对奠定一致的数据基础。

\subsubsection{（3）可视化智能系统：全流程闭环应用开发}

\begin{figure}[htbp]
    \centering
    \fbox{\parbox[c][4.5cm][c]{0.8\textwidth}{\centering 【图片空位：image1.png】黑土地智能监测系统核心功能界面展示}}
    \caption{黑土地智能监测系统核心功能界面展示}
    \label{fig:1}
\end{figure}

开发并测试了一套可视化智能系统，集成了从数据处理到结果输出的全流程功能：

\begin{itemize}
    \item 支持单张/批量遥感影像变化检测；
    \item 提供结果可视化展示、AI智能解译与一键生成分析报告的功能；
    \item 实现从数据输入到决策支撑的自动化流程，降低用户使用门槛，提升监测效率。
\end{itemize}

系统采用前后端分离的轻量化架构：后端基于FastAPI构建，负责影像接收、任务调度、模型推理与结果组织，接口结构清晰、便于维护与扩展；前端为可视化Web系统，围绕“影像上传—变化检测—结果可视化—AI解读—一键报告”组织操作流程，用户无需掌握遥感专业知识即可完成一次完整的监测作业。批量检测功能解决了连续时相、多个图斑的重复操作问题，使模型的批量推理能力能够真正转化为业务效率；AI解读环节把模型输出的变化图斑转述为自然语言结论，帮助基层人员理解变化的位置、类型与可能成因；一键报告则将影像、图斑与结论自动汇总为规范文档，便于上报归档与横向比对。上述设计使系统从单一的算法工具转变为面向决策的完整应用，普通电脑即可本地化部署，为黑土地保护的常态化监测与精细化治理提供了可持续运转的技术支撑。
